\section{From split point kernels to the remaining spanning problem}

% BEGIN J74T1
\paragraph{J74T1. The one-class normalization for an arbitrary split move.}
Fix the original pinned base record and a finite split move $h$ off its
bad integer as in S74B4. Write $L_h$ for its unconjugated lattice,
$S=\mathcal O_E^n$, $I=[L_h:L_h\cap S]$ and $r=D.R.r$.
For the geometric lattice $\Lambda=\bar L_h$ let $m$ and $Z$
denote $\mathrm{mIdx}(\Lambda,N,r)$ and
$\mathrm{zCard}(\Lambda,N,r)$ at an allowed denominator $N$.
Let $\chi(y)$ be the original lift character on the supported
determinant root $y$. Then
\[
        \sigma\bigl(\mathrm{epF}(L_h)Z/m\bigr)
                       =\chi(y)^{-1}\sqrt{I^r}.
\]
Here $\sigma$ embeds the coefficient field of the base record. The
identity applies to a move produced after the raw short series was
chosen; $h$ need not be one of the representatives originally stored
in the class sum.

Indeed, put $d_P=v_P(\det h_f)$. Unitarity gives
$d_{\bar P}=-d_P$ and off the moved split pairs $d_P=0$. Define
$\mathfrak c$ by the exponents $\max(d_P,0)$. The support and
coprimality hypotheses give
$\mathfrak c+\bar{\mathfrak c}=\mathcal O_E$ and
$\det(L_h)\bar{\mathfrak c}=\mathfrak c$. At a moved pair the
unramified local character values on uniformizers are inverse since
the character restricts trivially to the diagonal split base field
there. A negative supported determinant valuation therefore gives the
positive exponent at the conjugate prime in $\mathfrak c$; a positive
one gives that exponent at the original prime. At the other places
the supported determinant is one. This proves that the ideal character
on $\mathfrak c$ has value $\chi(y)$. By the two pins of RTG8--RTG9,
for $Q=\mathrm N(\mathfrak c)$ there is a \emph{positive} rational
$q$ with
\[
 q^2 I^r=Q,\qquad
 \sigma(\mathrm{epF}(L_h))=q/(\chi(y)\sqrt Q).
\]
The level guard puts $NS\subseteq\Lambda\cap S\subseteq\Lambda$.
Multiplicativity of indices in this chain gives $Z=mI^r$:
conjugation preserves $S$ and the relevant index. All indices are
positive and finite since the lattices are full and $N\ne0$.
Thus the left side equals $\chi(y)^{-1}q I^r/\sqrt Q$,
which is the claimed expression since $\sqrt Q=q\sqrt{I^r}$.
% END J74T1

% BEGIN J74T2
\paragraph{J74T2. The two point comparisons use one common multiplier.}
Work at a pinned base record with the core, frame and oscillator
extension hypotheses of P74I2--P74I7. Retain the original prime
agreement, family laws, and inclusion of the embedded coefficient
field in $O.c'.F$ required by T73C5 for the chosen record.
Take a normalization-one
short value $t(g)$ after the move of S74B4, at the dictionary frame
of an arbitrary target adelic point $g$. The harmonic pin represents
its polynomial by $P=O.\mathrm{io}(v)$.
Use the indicator $\phi$ of S74B1. We can now supply both point
comparisons required by S74B5 on these data.

To see the normalization geometrically, use the marked period chart
of the given frame. Write $\mathcal Q=(\mathrm{LevQ})^n$,
$\kappa'$ for the generator-to-vacuum comparison in this frame,
and $H$ for the evaluation of the classical derivative polynomial
on the projected vacuum translated by the dual roots and weighted
by $\widehat\phi$. In the complex model the normalization-one
short value is
\[
            t(g)=\frac{Z}{m|\mathcal Q|}\,\kappa' H .
\]
Here is how its three finite factors arise. Use representatives for
the finite isotropic lattice quotient, using the same even lifts
and the same marked roots as in P74I6 on the
complex fiber. They provide an admissible presentation in the
defining short-value rule; uniqueness in T73C1 permits this choice.
The canonical torsion lifts read as Heisenberg translates of the
theta generator. In the pulled-back chart the scaled variable is
$Nz$, so the homogeneous derivative of order $d$ acquires $N^d$
by the chain rule. This cancels its defining $N^{-d}$ factor.
The harmonic projection removes the quadratic frame correction under
differentiation, as in T73C4. The derivative is linear in the section,
so it commutes with both
finite sums. The sum of translated vacua over the isotropic quotient
is $Z$ times the averaging idempotent applied to the vacuum. The residue sum gives the
weighted translate of that projection, with the factor
$1/|\mathcal Q|$. The remaining normalization in the geometric
rule is $1/m$, and reading the product generator supplies $\kappa'$.
These identifications use the chart's actual basis, marked roots,
theta connection and renamed polynomial; changing the marking in
just one of the two sums would change $H$.

Write $C_{k(g)}$ for the component constant of P74I5 and
$K_d=a\delta^d(-1)^d(\sqrt{2\pi\xi_R})^{-d}$ as in P74I7,
now using the coefficient phase $a$ of this one base record.
The rational element moving $g$ into the marked chart preserves the
adelic theta evaluation. The base twist is trivial. P74I2 and P74I6
give for the point kernel
\[
 K_{h,v}(g)=K_d C_{k(g)}\mu(g)\,
       \chi(y)^{-1}\sqrt{I^r}\,\kappa'H/|\mathcal Q|.
\]
The residue Fourier weight is the embedded geometric weight by RTG11.
Substitution of J74T1 gives
\[
 K_{h,v}(g)=K_d\, C_{k(g)}\mu(g)\,
                   \sigma(\mathrm{epF}(L_h))\, t(g).
\]
No geometric value has been divided out.

For an integrated lift $\Theta_j$ of a rational source vector take
its placed class coefficients and the section $f_j$ of T73C5.
Use this \emph{same} choice of $a$, sign, component constants and
charts in the finite sum of P74I7. Its calculation on a marked
adelic point precedes passage to the quotient representative, so it
applies here at each $g$. The pinned cover constant $c_0$ of this
base record is a root of unity. It gives
\[
 c_0\Theta_j(g)=K_d C_{k(g)}\mu(g)\,
                         \operatorname{rd}_{\mathrm{dict}} f_j(g).
\]
In both equations $C_{k(g)}\mu(g)$ is the same nonzero function.
In the notation of S74B5, $\ell=(\sqrt{2\pi\xi_R})^{-d}$,
so the required scalar for the integrated lifts is
$u=K_d/(c_0\ell)$, of modulus one. If
$K_{h,v}=\sum b_j\Theta_j$ as in K73G3, the equations yield the
somewhat simpler identity
\[
       t=\frac{1}{c_0\,\sigma(\mathrm{epF}(L_h))}
                 \sum_j b_j\operatorname{rd}_{\mathrm{dict}} f_j .
\]
The denominator is nonzero by J74T1. This proves generation for
every short value reached by a split move, including a value which
vanishes identically.

Consequently the stable identity-cover expansion and the analytic
change-of-vacuum formula retained in S74B2--S74B4 suffice for the
original $\mathrm{ClassSumsSpan}$ on this pinned base record: those
arguments bring the cusp readings into the span of the moved short
values, and S74B5 descends that span to its coefficient field.
The higher covers on the prescribed field are available under D74U5,
including the section for the trivial determinant character. T73C5
supplies class sums also after the allowed change of phase.
P74I7 supplies the common point identity. RTG16--RTG18 can therefore
transfer the resulting base span to the original all-tuple comparison
under their stated base, phase and coefficient hypotheses. The
analytic expansion and change-of-vacuum formula, and the outer
construction of one family with the required inputs, remain separate
existence requirements.
% END J74T2
